% GENERATED FILE. Edit canonical lessons, then run npm run generate:books.
% canonical-source-hash: fnv1a64:86d72446a3abeb68
% canonical-lessons: GE-C217-kaputt, GE-C217-klug, GE-C217-kostenlos, GE-C217-kraeftig, GE-C217-ledig

\chapter{Broken, Clever, Free of charge, Strong, Single}
\label{ch:ge-broken-clever-free-of-charge-strong-single}
\hlchaptermodality{217}

\begin{chapteropening}
\textbf{By the end of this chapter:} \emph{I can say broken, clever, free of charge, strong and single.}

Complete the last lesson of chapter 217: I can say broken, clever, free of charge, strong and single.
\end{chapteropening}

\section[kaputt]{kaputt — broken}

\label{lesson:GE-C217-kaputt}



\begin{quote}
Before the new one: say the German for helpful, then the German for intelligent.
\end{quote}

\subsection*{You'll want to know: kaputt}
\textbf{kaputt} — \textquotedblleft{}broken\textquotedblright{}.

\begin{cousinweb}
\textbf{kaputt} $\leftarrow$ French \emph{capot}, said of a card player who has won no trick (EtymDB 2.1) — borrowed: the loser was done for. English \emph{kaput} came from German.
\end{cousinweb}

\begin{grammarlens}[title={What you've built}]
One more word for this chapter.
\end{grammarlens}

\subsection*{Your turn}
\begin{itemize}
\raggedright
  \item \emph{Say it:} \emph{kaputt}
  \item \emph{Say it:} \emph{kaputt}, once more
  \item \emph{Say it:} say \emph{kaputt}
  \item \emph{Recall:} say the German for helpful, then the German for intelligent, then say \emph{kaputt} again
\end{itemize}

\begin{tcolorbox}[breakable,colback=teal!4,colframe=teal!35!black,title={Before you move on}]
What does kaputt mean? (\textquotedblleft{}Broken\textquotedblright{}.) Say it once more.
\end{tcolorbox}
\section[klug]{klug — clever}

\label{lesson:GE-C217-klug}



\begin{quote}
Before the new one: say the German for intelligent, then the German for broken.
\end{quote}

\subsection*{You'll want to know: klug}
\textbf{klug} — \textquotedblleft{}clever\textquotedblright{}.

\begin{cousinweb}
\textbf{klug} $\leftarrow$ Middle High German \emph{kluoc}, \textquotedblleft{}clever, fine\textquotedblright{} (EtymDB 2.1) — inherited; EtymDB lists Old Norse \emph{klókr}, \textquotedblleft{}clever\textquotedblright{}, beside it.
\end{cousinweb}

\begin{grammarlens}[title={What you've built}]
One more word for this chapter.
\end{grammarlens}

\subsection*{Your turn}
\begin{itemize}
\raggedright
  \item \emph{Say it:} \emph{klug}
  \item \emph{Say it:} \emph{klug}, once more
  \item \emph{Say it:} say \emph{klug}
  \item \emph{Recall:} say the German for intelligent, then the German for broken, then say \emph{klug} again
\end{itemize}

\begin{tcolorbox}[breakable,colback=teal!4,colframe=teal!35!black,title={Before you move on}]
What does klug mean? (\textquotedblleft{}Clever\textquotedblright{}.) Say it once more.
\end{tcolorbox}
\section[kostenlos]{kostenlos — free of charge}

\label{lesson:GE-C217-kostenlos}



\begin{quote}
Before the new one: say the German for broken, then the German for clever.
\end{quote}

\subsection*{You'll want to know: kostenlos}
\textbf{kostenlos} — \textquotedblleft{}free of charge\textquotedblright{}.

\begin{cousinweb}
\textbf{kostenlos} is the word for \textquotedblleft{}costs\textquotedblright{}, from \emph{kosten}, + \emph{-los}, \textquotedblleft{}without\textquotedblright{}, the same ending as English \emph{-less}: without cost. EtymDB 2.1 lists English \emph{costless} beside it.
\end{cousinweb}

\begin{grammarlens}[title={What you've built}]
One more word for this chapter.
\end{grammarlens}

\subsection*{Your turn}
\begin{itemize}
\raggedright
  \item \emph{Say it:} \emph{kostenlos}
  \item \emph{Say it:} \emph{kostenlos}, once more
  \item \emph{Say it:} say \emph{kostenlos}
  \item \emph{Recall:} say the German for broken, then the German for clever, then say \emph{kostenlos} again
\end{itemize}

\begin{tcolorbox}[breakable,colback=teal!4,colframe=teal!35!black,title={Before you move on}]
What does kostenlos mean? (\textquotedblleft{}Free of charge\textquotedblright{}.) Say it once more.
\end{tcolorbox}
\section[kräftig]{kräftig — strong, hearty}

\label{lesson:GE-C217-kraeftig}



\begin{quote}
Before the new one: say the German for clever, then the German for free of charge.
\end{quote}

\subsection*{You'll want to know: kräftig}
\textbf{kräftig} — \textquotedblleft{}strong, hearty\textquotedblright{}.

\begin{cousinweb}
\textbf{kräftig} is the word for \textquotedblleft{}strength\textquotedblright{}, Old High German \emph{kraft} $\leftarrow$ Proto-Germanic \emph{*kraftaz} (EtymDB 2.1), + \emph{-ig}. English \emph{craft}, first \textquotedblleft{}strength, skill\textquotedblright{}, is the same word.
\end{cousinweb}

\begin{grammarlens}[title={What you've built}]
One more word for this chapter.
\end{grammarlens}

\subsection*{Your turn}
\begin{itemize}
\raggedright
  \item \emph{Say it:} \emph{kräftig}
  \item \emph{Say it:} \emph{kräftig}, once more
  \item \emph{Say it:} say \emph{kräftig}
  \item \emph{Recall:} say the German for clever, then the German for free of charge, then say \emph{kräftig} again
\end{itemize}

\begin{tcolorbox}[breakable,colback=teal!4,colframe=teal!35!black,title={Before you move on}]
What does kräftig mean? (\textquotedblleft{}Strong, hearty\textquotedblright{}.) Say it once more.
\end{tcolorbox}
\section[ledig]{ledig — single (unmarried)}

\label{lesson:GE-C217-ledig}



\begin{quote}
Before the new one: say the German for free of charge, then the German for strong.
\end{quote}

\subsection*{You'll want to know: ledig}
\textbf{ledig} — \textquotedblleft{}single (unmarried)\textquotedblright{}.

\begin{cousinweb}
\textbf{ledig} $\leftarrow$ Old High German \emph{lidig}, \textquotedblleft{}free\textquotedblright{} $\leftarrow$ Proto-Germanic \emph{*liþugaz}, \textquotedblleft{}flexible, free, unoccupied\textquotedblright{} — inherited. (EtymDB 2.1) Single is free of ties.
\end{cousinweb}

\begin{grammarlens}[title={What you've built}]
One more word for this chapter.
\end{grammarlens}

\subsection*{Your turn}
\begin{itemize}
\raggedright
  \item \emph{Say it:} \emph{ledig}
  \item \emph{Say it:} \emph{ledig}, once more
  \item \emph{Say it:} say \emph{ledig}
  \item \emph{Recall:} say the German for free of charge, then the German for strong, then say \emph{ledig} again
\end{itemize}

\begin{tcolorbox}[breakable,colback=teal!4,colframe=teal!35!black,title={Before you move on}]
What does ledig mean? (\textquotedblleft{}Single (unmarried)\textquotedblright{}.) Say it once more.
\end{tcolorbox}
